\documentclass[11pt]{memoir}
\usepackage{raeez-math-template}

\title{Mathematical Typography Specimen}
\author{Raeez Lorgat}
\date{\today}

\begin{document}
\maketitle

\begin{abstract}
This file is the executable specimen for the shared mathematics template.
It exercises the page colour, small-cap hierarchy, theorem surface, display
spacing, tables, citations, and dense mathematical notation.
\end{abstract}

\tableofcontents

\chapter{Optical System}

\section{Text Colour and Measure}

The template is tuned for long mathematical manuscripts: a quiet Garamond
text face, restrained line length, old-style proportional figures in prose,
and a math face whose italic, subscripts, delimiters, and operators are
matched to the surrounding paragraph.  The page should feel calm before it
feels ornate.  Conceptual hierarchy is carried by theorem names, labels, and
proof structure; visual hierarchy remains deliberately narrow.

\begin{definition}[Garamond mathematical page]
A mathematical page is \emph{balanced} if the text block, display equations,
theorem headings, footnotes, and tables share one optical weight and one
baseline rhythm.
\end{definition}

\begin{theorem}[Display rhythm]
Let $X$ be a compact Calabi--Yau threefold and let
$\mathcal C = \operatorname{Perf}(X)$.  If
$\operatorname{HH}^{\bullet}(\mathcal C)$ carries its cyclic
$E_3$-structure, then the schematic chiral envelope
\[
  \Phi_{\mathrm{ch}}(\mathcal C)
  \simeq
  \int_{S^2}
  \operatorname{Obs}_{\mathrm{hCS}}(X)
  \otimes_{\mathbb Q}
  \mathbb Q[[\hbar]]
\]
has a typographic presentation whose displayed operators do not dominate
the paragraph that proves them.
\end{theorem}

\begin{proof}
The typographic content is not the formula itself, but the relation among
the paragraph, the display, and the theorem head.  The package fixes the
display skip, operator scale, Garamond math metrics, and small-cap tracking
at the template level, leaving the manuscript to carry only the mathematics.
\end{proof}

\section{Dense Formulae}

The following display tests operators, large delimiters, blackboard bold,
script letters, and nested subscripts:
\[
  \mathfrak g_{\Delta_5}
  =
  \bigoplus_{\alpha \in \Lambda^{2,1}_{II,+}}
  \mathfrak g_{\alpha},
  \qquad
  \sum_{n,m \geq 0}
  \operatorname{sdim}
  \mathfrak g_{(n,\ell,m)}
  q^n r^\ell s^m
  =
  \log\!\left(
    \prod_{\gamma \in \Gamma_{\mathrm{eff}}}
    (1-x^\gamma)^{-c(\gamma)}
  \right).
\]

\begin{remark}
Tables use the same visual economy: booktabs rules, modest row leading, and
no vertical boxing.
\end{remark}

\section{Kerning and Symbol Pressure}

The pairwise texture matters in running prose: AV, WA, Ta, Te, To, Yo, Ly,
``quoted phrases,'' parenthetical clauses, and em-dashes---not heavy rules in
miniature---must all disappear into the sentence.  Figures such as 1729,
2026, and 3.14159 use old-style proportions in prose, while labels and
tables keep their alignment when the manuscript asks for it.\footnote{The
footnote is intentionally ordinary: it tests smaller Garamond colour, the
footnote rule, old-style figures, and the distance between note text and the
main block.}

Inline mathematics must keep the same colour: $\mathbb Q \subset \mathbb C$,
$\mathcal F_{\hbar}^{\mathrm{ch}}$, $\mathfrak g_{\Delta_5}$,
$\mathscr L_{\log}$, $\operatorname{RHom}_{\mathcal A}(M,N)$, and
$E_3 \otimes \mathbb Q[[\hbar]]$ should read as mathematical prose rather
than as dropped-in symbols.

\begin{align}
  \operatorname{RHom}_{\mathcal A}(M,N)
  &\simeq
  \operatorname{Tot}
  \left[
    \prod_{p+q=n}
    \operatorname{Hom}_{\mathbb Q}
    \bigl(B_p\mathcal A \otimes M_q,N\bigr)
  \right],\\
  \Delta_5(Z)
  &=
  e^{2\pi i(\rho,Z)}
  \prod_{\alpha\in\mathcal R_+}
  \left(1-e^{-2\pi i(\alpha,Z)}\right)^{\operatorname{mult}(\alpha)}.
\end{align}

\begin{center}
\begin{tabular}{lll}
\toprule
Layer & Typographic responsibility & Mathematical responsibility\\
\midrule
Text & measure, colour, kerning & exposition\\
Display & rhythm, delimiter weight & identities\\
Theorem & hierarchy, head weight & logical dependency\\
\bottomrule
\end{tabular}
\end{center}

\begin{biblist}
\bib{knuth}{book}{
  author={Knuth, Donald E.},
  title={The TeXbook},
  publisher={Addison-Wesley},
  date={1984},
}
\end{biblist}

\end{document}
